\section{Geometric Primitives}

The motion-planning system begins with a small collection of geometric
primitives. These structures provide the representation used by all later
geometry and planning algorithms.

\subsection{Point}

A two-dimensional point is represented as
\[
p=(x,y), \qquad x,y\in\mathbb{R}.
\]

In the implementation, coordinates are stored using double-precision
floating-point values. A point is represented in C++ as:

\begin{lstlisting}[language=C++]
struct Point {
    double x;
    double y;
};
\end{lstlisting}

\subsection{Line Segment}

A line segment is defined by two endpoints $A$ and $B$ and is denoted by
$\overline{AB}$. If
\[
A=(x_A,y_A), \qquad B=(x_B,y_B),
\]
then the segment contains every point
\[
P(t)=A+t(B-A), \qquad 0\leq t\leq 1.
\]

The C++ representation is:

\begin{lstlisting}[language=C++]
struct Segment {
    Point a;
    Point b;
};
\end{lstlisting}

\subsection{Polygon}

A polygon is represented as an ordered sequence of vertices
\[
\mathcal{P}=\{p_0,p_1,\ldots,p_{n-1}\}.
\]

Each consecutive pair of vertices forms an edge, and the final vertex
connects back to the first. Thus the edges are
\[
\left(p_i,p_{(i+1)\bmod n}\right), \qquad 0\leq i<n.
\]

This modular indexing automatically generates the closing edge
$p_{n-1}\rightarrow p_0$.

For this project, a valid polygon contains at least three vertices. Polygons
with fewer than three vertices are treated as invalid and therefore produce no
polygon edges.

Obstacles are assumed to be \emph{simple polygons} unless otherwise stated. A
simple polygon has no self-intersections between non-adjacent edges.
Self-intersecting polygons are outside the current geometric model.

\subsection{Path}

A robot path is represented using an ordered list of waypoints
\[
P_0,P_1,\ldots,P_k.
\]

The corresponding path segments are
\[
(P_0,P_1),(P_1,P_2),\ldots,(P_{k-1},P_k).
\]

Storing waypoints rather than explicitly storing every segment avoids
duplicating information and simplifies path modification.
